\par\medskip\noindent
\underline{Lema 3.5.} Considérese un diagrama de $R$-álgebras
\[
\begin{tikzcd}[column sep=large,row sep=large]
 & C \arrow[d]\\
 B' \arrow[r,"\text{suprayectiva}"'] & B
\end{tikzcd}
\]
y sea $S$ una $R$-álgebra.

\begin{enumerate}[label=(\roman*)]
\item Si $\operatorname{Tor}_1^R(B,S)=0$, entonces la aplicación canónica
\[
 (B'\times_BC)\otimes_RS
 \longrightarrow
 (B'\otimes_RS)\times_{B\otimes_RS}(C\otimes_RS)
\]
es un isomorfismo.

\item Supongamos que $B'$ es plano sobre $R$. Si
\[
 \operatorname{Tor}_i^R(C,S)=0\quad(0<i<n)
 \qquad\text{y}\qquad
 \operatorname{Tor}_i^R(B,S)=0\quad(0<i\sisi{<}{\leq}n),
\]
entonces
\[
 \operatorname{Tor}_i^R(B'\times_BC,S)=0\qquad(0<i<n).
\]
\end{enumerate}

\underline{Demostración.} Considérese la sucesión exacta
$0\longrightarrow J\longrightarrow B'\longrightarrow B
\longrightarrow0$.

(i) Si $\operatorname{Tor}_{\sisi{i}{1}}^R(B,S)=0$, entonces
\[
 0\longrightarrow J\otimes_RS\longrightarrow B'\otimes_RS
 \longrightarrow B\otimes_RS\longrightarrow0
\]
es exacta, y por tanto también lo es
\[
 0\longrightarrow J\otimes_RS
 \longrightarrow
 (B'\otimes_RS)\times_{B\otimes_RS}(C\otimes_RS)
 \longrightarrow C\otimes_RS\longrightarrow0.
\]
Se obtiene así un diagrama conmutativo de filas exactas:
\begin{center}
\scriptsize
\begin{tikzcd}[column sep=small,row sep=large]
 & J\otimes_RS \arrow[r] \arrow[d,equal]
 & (B'\times_BC)\otimes_RS \arrow[r] \arrow[d]
 & C\otimes_RS \arrow[r] \arrow[d,equal]
 & 0\\
 0 \arrow[r] & J\otimes_RS \arrow[r]
 & (B'\otimes_RS)\times_{B\otimes_RS}(C\otimes_RS) \arrow[r]
 & C\otimes_RS \arrow[r] & 0.
\end{tikzcd}
\end{center}

%%%%%%%%%%  scan idx 88  |  folio 77  %%%%%%%%%%

Se sigue que
\[
 (B'\sisi{\otimes}{\times}_BC)\otimes_RS
 \longrightarrow
 (B'\otimes_RS)\times_{B\otimes_RS}(C\otimes_RS)
\]
es un isomorfismo.

(ii) Puesto que $B'$ es plano sobre $R$ y
$\operatorname{Tor}_i^R(B,S)=0$ para $0<i\leq n$, la sucesión exacta
larga de Tor da
\[
 \operatorname{Tor}_i^R(J,S)=0\qquad(0<i<n).
\]
Aplicando la misma sucesión a
\[
 0\longrightarrow J\longrightarrow B'\times_BC
 \longrightarrow C\longrightarrow0
\]
y usando la hipótesis sobre $C$, se obtiene
\[
 \operatorname{Tor}_i^R(B'\times_BC,S)=0\qquad(0<i<n).
\]
\hfill$\square$

\underline{Corolario 3.6.}

\begin{enumerate}[label=(\alph*)]
\item Sean $A$ y $S$ $R$-álgebras. Si
$\operatorname{Tor}_1^R(A,S)=0$, entonces, para todo
$G\in\operatorname{Ob}
(\overset{\vee}{\underline C}_{A\otimes_RS|S})$, la aplicación canónica
\[
 \check H^{\bullet}(A\otimes_RS|S,G)
 \longrightarrow\check H^{\bullet}(A|R,\check SG)
\]
es un isomorfismo. Aquí
\[
 S:\underline C_{A|R}\longrightarrow
 \underline C_{A\otimes_RS|S},
 \qquad X\longmapsto X\otimes_RS.
\]
En particular, si $S$ es plano sobre $R$, entonces
\[
 \check S\,\check{\underline H}^{\,\bullet}(G)
 \longrightarrow
 \check{\underline H}^{\,\bullet}(\check SG)
\]
es un isomorfismo para todo $G$ de esta clase.

\item Sea $P\longrightarrow B$ una aplicación suprayectiva de
$R$-álgebras, donde $P$ es un álgebra polinómica sobre $R$. Si
\[
 \operatorname{Tor}_i^R(B,S)=0\qquad(0<i\leq n),
\]
entonces, para todo producto fibrado iterado de copias de $P$ sobre $B$,
\[
 \operatorname{Tor}_i^R
 (P\times_B P\times_B\cdots\times_B P,S)=0
 \qquad(0<i<n).
\]
\end{enumerate}

\underline{Demostración.}

(a) Dado $B\in\operatorname{Ob}(\underline C_{A|R})$, elíjase un
recubrimiento $\{P\longrightarrow B\}\in\operatorname{Cov}(B)$ con $P$
un álgebra polinómica sobre $R$. Si
$\operatorname{Tor}_1^R(B,S)=0$, el Lema 3.5(i) muestra que la aplicación
natural
\[
 \left(\overbrace{P\times_B\cdots\times_BP}^{n}\right)\otimes_RS
 \longrightarrow
 \overbrace{(P\otimes_RS)\times_{B\otimes_RS}\cdots
 \times_{B\otimes_RS}(P\otimes_RS)}^{n}
\]
es un isomorfismo para todo $n$. Por consiguiente,
\begin{align*}
 \check H^{\bullet}(B|R,\check SG)
 &=H^{\bullet}(\{P\longrightarrow B\},\check SG)\\
 &=H^{\bullet}(\{P\otimes_RS\longrightarrow B\otimes_RS\},G)\\
 &=\check H^{\bullet}(B\otimes_RS|S,G).
\end{align*}
Si $S$ es plano sobre $R$, entonces
$\operatorname{Tor}_1^R(X,S)=0$ para todo
$X\in\operatorname{Ob}(\underline C_{A|R})$. Por tanto, el isomorfismo
precedente es natural en $X$, y da
\[
 \check S\,\check{\underline H}^{\,\bullet}(G)
 \xrightarrow{\ \sim\ }
 \check{\underline H}^{\,\bullet}(\check SG).
\]

(b) Pongamos
\[
 P^{(m)}=\overbrace{P\times_B\cdots\times_BP}^{m+1}.
\]
Para $m=0$ no hay nada que demostrar, pues $P$ es plano sobre $R$.
Supongamos inductivamente que
$\operatorname{Tor}_i^R(P^{(m-1)},S)=0$ para $0<i<n$. Puesto que
$P\longrightarrow B$ es suprayectiva y
$P^{(m)}=P\times_BP^{(m-1)}$, el Lema 3.5(ii) da
\[
 \operatorname{Tor}_i^R(P^{(m)},S)=0\qquad(0<i<n).
\]
\hfill$\square$
